\section{Results}\label{sec:results}

\subsection{Validation}\label{sec:results-validation}
All checks of Section~\ref{sec:method-validation} passed before the measurements. The MILP certified the optimum of all 30 instances of E1, the optimum over the enumerated space $\Fk$ matched the MILP on all 30, and all instances were non-degenerate ($0<\Fstar<F(0)$). An independent maximum flow computation confirmed the optimality of each of the 18{,}000 runs that reported an optimal plan.

\subsection{Scaling of the cost to the optimum (H1)}\label{sec:results-scaling}
The mean cost of minimum finding grew with the square root of the number of feasible plans, and that of random-order search grew linearly (Fig.~\ref{fig:scaling}). Against $N$, the fitted exponent of minimum finding was $\beta=0.494$ (95\% CI $[0.460,0.523]$, $R^2=0.977$), inside the registered interval $[0.4,0.6]$, so H1 is not refuted. The confidence interval also contains the theoretical value $0.5$. Random-order search had $\beta=0.975$ ($[0.940,1.006]$). Against $N/M$, which accounts for the number of optimal plans, the exponents were $0.510$ ($[0.497,0.529]$, $R^2=0.995$) and $1.001$ ($[0.984,1.022]$), respectively.

On the three largest instances ($N=974{,}121$), minimum finding reached the optimum after 1{,}750 queries and evaluations on average, against 383{,}687 evaluations for random-order search and 974{,}121 for enumeration. Across instances, the mean cost of minimum finding ranged from $0.90\sqrt{N}$ to $2.67\sqrt{N}$, and the lowest values occurred on instances with several optimal plans ($M$ up to 4). All 6{,}000 runs of minimum finding reached the optimum within the budget $Q$ of Proposition~\ref{prop:advantage}, counting queries and evaluations, and the mean cost never exceeded 10.1\% of $Q$. Averaged over the instances, the mean cost of random-order search was within 1\% of the theoretical value $(N+1)/(M+1)$.

\begin{figure}[t]
\centering
\begin{tikzpicture}[font=\footnotesize]
\begin{loglogaxis}[
  width=\linewidth,
  height=5.3cm,
  xmin=5, xmax=2.5e6,
  ymin=1, ymax=2.5e6,
  xtickten={1,...,6},
  ytickten={0,...,6},
  xminorticks=false,
  yminorticks=false,
  tick pos=left,
  tick align=outside,
  major tick length=2.5pt,
  xlabel={Number of feasible plans $N$},
  ylabel={Mean cost to the optimum},
  label style={font=\footnotesize},
  tick label style={font=\footnotesize},
  axis line style={line width=0.4pt},
  every tick/.style={black, line width=0.4pt},
  table/col sep=comma,
]
\addplot[qgray, line width=0.6pt, dotted] table[x=N, y=N] {figures/data/scaling_fits.csv}
  node[pos=1, anchor=west, text=qgray] {$N$};
\addplot[black, line width=0.6pt] table[x=N, y=random_order] {figures/data/scaling_fits.csv}
  node[pos=0.8, sloped, anchor=north, yshift=-2pt] {random-order search};
\addplot[qaccent, line width=0.6pt] table[x=N, y=durr_hoyer] {figures/data/scaling_fits.csv}
  node[pos=0.8, sloped, anchor=north, yshift=-3pt, text=qaccent] {minimum finding};
\addplot[only marks, mark=square, mark size=1.2pt, mark options={black, line width=0.5pt, fill=white}] table[x=N, y=random_order] {figures/data/scaling_points.csv};
\addplot[only marks, mark=*, mark size=1.2pt, mark options={qaccent, line width=0.5pt, fill=qaccent}] table[x=N, y=durr_hoyer] {figures/data/scaling_points.csv};
\end{loglogaxis}
\end{tikzpicture}

\caption{Mean cost to the optimum of minimum finding and random-order search on the 30 instances of E1 (200 runs per instance), with least-squares fits and the enumeration cost $N$ for reference.}\label{fig:scaling}
\end{figure}

\subsection{Resources of the interface}\label{sec:results-resources}
The Dicke preparation was the cheapest component and the arithmetic oracle the most expensive (Table~\ref{tab:resources}). For networks with 3 to 14 arcs, the preparation $\Ak$ used 8 to 174 two-qubit gates in the logical count. The truth-table oracle grew quickly with $m$ and with the number of marked plans. The arithmetic oracle over $(x,\alpha)$ used 18 to 46 qubits, of which only $m+n-2$ take part in the search, while the others are ancillas for the cut indicators, the weighted sum, and the comparison. It needed 690 to 4{,}202 two-qubit gates in the logical count and 2{,}055 to 14{,}296 after transpilation to the Heron model, where routing on the heavy-hex connectivity roughly triples the count. The condition $2^{n-2}<N$ held on 7 of the 9 instances and failed on the two instances with $k=1$ and $n\ge 6$.

\begin{table*}[t]
\centering
\caption{Two-qubit gate counts of the interface components on the instances of E2, logical (basis \texttt{cx}, \texttt{u}) / transpiled for \texttt{FakeFez} at optimization level 1. The truth-table oracle marks the $M$ optimal plans, and the last column states whether $2^{n-2}<N$.}\label{tab:resources}
\small
\begin{tabular}{rrrrrrrrc}
\toprule
$n$ & $m$ & $k$ & $N$ & $A_k$ & Truth table & Arith. qubits & Arith. oracle & $2^{n-2}<N$ \\
\midrule
3 & 3 & 1 & 4 & 8 / 8 & 6 / 9 & 18 & 690 / 2{,}055 & yes \\
4 & 7 & 1 & 8 & 24 / 24 & 122 / 242 & 29 & 2{,}008 / 6{,}533 & yes \\
4 & 7 & 2 & 29 & 76 / 121 & 122 / 242 & 29 & 2{,}014 / 6{,}353 & yes \\
5 & 10 & 1 & 11 & 36 / 36 & 290 / 686 & 34 & 2{,}344 / 7{,}465 & yes \\
5 & 10 & 2 & 56 & 118 / 190 & 290 / 686 & 34 & 2{,}338 / 7{,}315 & yes \\
6 & 12 & 1 & 13 & 44 / 44 & 442 / 1{,}096 & 43 & 4{,}202 / 14{,}296 & no \\
6 & 12 & 2 & 79 & 146 / 236 & 442 / 1{,}096 & 43 & 4{,}202 / 14{,}296 & yes \\
7 & 14 & 1 & 15 & 52 / 52 & 1{,}252 / 3{,}058 & 46 & 4{,}140 / 13{,}615 & no \\
7 & 14 & 2 & 106 & 174 / 282 & 626 / 1{,}514 & 46 & 4{,}146 / 13{,}927 & yes \\
\bottomrule
\end{tabular}
\end{table*}


\subsection{Amplitude amplification on quantum hardware (H2)}\label{sec:results-hardware}
The E3 instance has three arcs, $s\to 1$, $s\to t$, and $1\to t$, with capacities 7, 10, and 6, maximum flow $F(0)=16$, and a single optimal plan, which removes $s\to t$ and leaves $\Fstar=6$. On the ideal simulator, the frequency of the optimal plan was within three standard errors of theory in all eight circuits. Table~\ref{tab:hardware} lists the frequencies at the optimal number of iterations. On the IBM Heron processor \texttt{ibm\_fez} (job \texttt{dbmujid8mdtc73d8m2h0}), the optimal plan was measured with frequency $0.815$ in the interface variant after one iteration and $0.818$ in the hypercube variant after two, well above the baselines $1/4$ and $1/8$ (one-sided binomial test, $p<10^{-300}$ in both cases), so H2 is not refuted. The noise model \texttt{FakeFez} gave $0.898$ and $0.833$ for the same circuits.

In the exploratory comparison, the hardware frequency in the interface variant (95\% Wilson interval $[0.803,0.827]$) reached 81.5\% of the theoretical probability of one and was significantly lower than on \texttt{FakeFez} (difference $0.083$, two-sided Fisher exact test, $p<10^{-26}$). In the hypercube variant, the hardware frequency ($[0.806,0.829]$) reached 86.5\% of the theoretical probability of $0.945$ and did not differ significantly from \texttt{FakeFez} (difference $0.016$, $p=0.06$). With the interface preparation, the frequency of the optimum dropped back to $0.215$ and $0.233$ after two and three iterations, close to the theoretical value of $0.25$. In the interface variant, the fraction of shots inside $\mathcal{F}_1$ fell from $0.983$ without iterations to $0.917$ after one iteration and $0.880$ after three. The job waited 41.2 minutes in the queue of the free plan and used 11 seconds of processor time.

\begin{table*}[t]
\centering
\caption{Frequency of the optimal plan at the optimal number of Grover iterations $j$ for the E3 instance (4{,}096 shots per circuit), with circuit depth and two-qubit gates after transpilation and the one-sided binomial $p$-value against $1/N$.}\label{tab:hardware}
\small
\begin{tabular}{llrrrrrrr}
\toprule
Environment & Variant & $j$ & Depth & 2q gates & Theory & Frequency & $1/N$ & $p$ \\
\midrule
ideal simulator & interface & 1 & -- & -- & 1.000 & 1.000 & 0.250 & $<10^{-300}$ \\
ideal simulator & hypercube & 2 & -- & -- & 0.945 & 0.942 & 0.125 & $<10^{-300}$ \\
noise model (\texttt{FakeFez}) & interface & 1 & 110 & 29 & 1.000 & 0.898 & 0.250 & $<10^{-300}$ \\
noise model (\texttt{FakeFez}) & hypercube & 2 & 144 & 39 & 0.945 & 0.833 & 0.125 & $<10^{-300}$ \\
hardware (\texttt{ibm\_fez}) & interface & 1 & 110 & 29 & 1.000 & 0.815 & 0.250 & $<10^{-300}$ \\
hardware (\texttt{ibm\_fez}) & hypercube & 2 & 144 & 39 & 0.945 & 0.818 & 0.125 & $<10^{-300}$ \\
\bottomrule
\end{tabular}
\end{table*}


\subsection{Classical context}\label{sec:results-classical}
The MILP solved every instance of E1 in at most 24\,ms, with no visible growth with $N$, whereas single-threaded enumeration of $\Fk$ grew linearly with $N$ and took up to 568\,s on the largest instances. In the MILP scaling study, with general costs and a limit of 20\,s, the MILP certified all three seeds of every size up to 2{,}560 nodes and 7{,}680 arcs in the random family with budgets of 50\% and 75\% of the cheapest cut. In the layered family it certified all seeds up to 4{,}098 nodes and 12{,}224 arcs, the largest size tested, with the same budgets. With a budget of 25\%, the corresponding sizes were 320 nodes in the random family, where one instance with 640 nodes exceeded the limit, and 2{,}306 nodes in the layered family.
